\section{The distinction that governs the problem}
A rational matrix, a rational Gram factor, and a rational point in a prescribed affine fiber are different objects. For a closed vector domain $P\subseteq\R^n$ that is a union of cones, write
\begin{align}
 C(P)&=\cone\{xx\trans:x\in P\},\\
 \GG(P)&=\left\{\sum_{j=1}^s u_ju_j\trans:s<\infty,\ u_j\in P\cap\Q^n\right\}.
\end{align}
The empty sum is permitted. Throughout the main theorem, $P$ is a finite union of closed rational polyhedral cones; components may have lineality or smaller dimension. In particular $P$ need not be convex. Let
\begin{equation}\label{eq:fiber}
 H=\{X\in\Ssym^n:\mathcal L(X)=b\}
\end{equation}
be rational affine, and put $d=n(n+1)/2$. For a convex cone $K$, $C(K)$ is the generalized completely positive cone and its dual consists of the quadratic forms nonnegative on $K$.

Our main object is a vector-domain invariant, not a pivot potential. Define the \emph{viable rational directions}
\begin{equation}\label{eq:viable}
 \mathcal D_{\Q}(P,H)=\{u\in P\cap\Q^n:\exists X\in\GG(P)\cap H,\ \exists\tau\in\Q_{>0},\ X-\tau uu\trans\in\GG(P)\}.
\end{equation}
The coefficient $\tau$ cannot be omitted: an arbitrarily long prescribed column cannot fit a fixed trace. Define
\begin{equation}\label{eq:carrierdef}
 \AC(P,H)=\cl\mathcal D_{\Q}(P,H).
\end{equation}
If the rational fiber is empty, its carrier is empty. If the only feasible matrix is zero, the carrier is $\{0\}$, not the empty set.

\begin{theorem}[Canonical arithmetic carrier]\label{thm:main}
For the preceding rational data there is a finite algorithm which either proves $\GG(P)\cap H=\varnothing$, or computes $\AC=\AC(P,H)$ as a finite union of closed rational polyhedral subcones of $P$. It uses at most $d$ proper matrix-span reductions. In the nonempty case,
\begin{align}
 \mathcal D_{\Q}(P,H)&=\AC\cap\Q^n,\label{eq:directions}\\
 \GG(P)\cap H&=\GG(\AC)\cap H,\label{eq:exactfiber}\\
 \cl(\GG(P)\cap H)&=C(\AC)\cap H,\label{eq:closure}\\
 H\cap\ri C(\AC)&\ne\varnothing.\label{eq:slater}
\end{align}
The set $\AC$, not merely the matrix slice in \eqref{eq:closure}, is independent of all choices of exposing forms. There is a rational $X_0\in H\cap\ri C(\AC)$ which witnesses every direction in \eqref{eq:directions}, with $\tau$ allowed to depend on $u$. Such a witness has a full-span reserve representation with at most $2d$ weighted rational atoms. Every feasible rational matrix admits at most $d$ weighted rational atoms, hence at most $4d$ rational Gram columns, whenever it belongs to $\GG(P)$.
\end{theorem}

The last width assertion is an existence and output-size statement; it is not a claim of a deterministic polynomial-time four-square routine. Section~\ref{sec:complexity} distinguishes that width from the polynomial-time output construction. The theorem does not assert that every rational point of the real fiber factors rationally. Rather, it computes the exact arithmetic obstruction.

Classical facial reduction supplies successive exposing functionals~\cite{BW,LMT}. The additional step here is at the level of \emph{column directions}: an algebraic zero restriction is replaced by its exact rational coefficient kernel before proceeding. Rational interior factorization~\cite{DSV17}, rational-point algorithms for convex semialgebraic sets~\cite{SEDZ}, and generalized perfect-matrix theory~\cite{OS26} are antecedents, not claims of this paper. In particular, deciding whether a rational matrix point exists in $C(P)\cap H$ does not decide whether that point has permitted rational columns.

The proof also explains a limitation of unification by a single facial depth. A perfect-form normal fan subdivides the \emph{interior} of a matrix cone. Its cells are not genuine faces of that cone. Section~\ref{sec:paths} constructs exponentially many consecutive cells while the genuine arithmetic carrier is unchanged. The corrected principle is therefore a certification theorem together with a precise obstruction to using it as a polynomial local-walk bound.

\section{Rank-one reserves and arithmetic zero domains}\label{sec:zeros}
We use the Frobenius pairing on symmetric matrices. A rational weighted representation means $X=\sum_j\lambda_j u_ju_j\trans$ with $\lambda_j\in\Q_{\ge0}$ and $u_j\in P\cap\Q^n$.

\begin{lemma}[Weights, closedness, and a rational reserve]\label{lem:reserve}
Rational weighted representations are equivalent to membership in $\GG(P)$. The real cone $C(P)$ is closed, its linear span $V$ is rational, and there are rational permitted atoms $G_i=a_ia_i\trans$, $1\le i\le k=\dim V$, forming a basis of $V$. Their sum $B=\sum_iG_i$ belongs to $\ri C(P)$.
\end{lemma}
\begin{proof}
For $\lambda=a/b>0$ in lowest terms, Lagrange's four-square theorem gives $ab=\sum_{\ell=1}^4s_\ell^2$ with nonnegative integers $s_\ell$. Thus
\[
 \lambda uu\trans=\sum_{\ell=1}^4((s_\ell/b)u)((s_\ell/b)u)\trans.
\]
Positive rescaling preserves every component cone. The reverse implication is immediate.

If $P\ne\{0\}$, the normalized atoms $xx\trans$ with $x\in P$, $\norm{x}_2=1$, form a compact set. Its convex hull is exactly the trace-one base of $C(P)$ and is compact. Since trace equals one on that base, its cone is closed. The zero case is immediate.

For each component choose rational cone vectors $a_1,\ldots,a_r$ spanning its vector span. Their squares and the squares of $a_i+a_j$ span the symmetric matrices supported on that span. All these vectors remain in the component. Taking the union over components and selecting an atom basis gives $G_1,\ldots,G_k$. The simplicial cone generated by that basis is full-dimensional in $V$; its positive sum has a relative neighborhood in that cone. Hence $B\in\ri C(P)$.
\end{proof}

\begin{lemma}[Rational relative-interior factorization]\label{lem:interior}
Every rational $X\in\ri C(P)$ belongs to $\GG(P)$. It has a rational weighted representation using at most $2k$ atoms that includes a strictly positive coefficient on every atom of a basis of $V$. Without the reserve requirement, at most $k$ weighted atoms suffice.
\end{lemma}
\begin{proof}
Choose rational $\tau>0$ sufficiently small that $X-\tau B\in C(P)$. Conic Carath\'eodory gives
\[
 X-\tau B=\sum_{j=1}^{k}x_jx_j\trans,
\]
allowing zero summands and absorbing real nonnegative weights into the vectors. Approximate each $x_j$ by rational $r_j$ in the same rational polyhedral component. In the rational basis $(G_i)$ the residual
\[
 X-\sum_jr_jr_j\trans=\sum_i e_iG_i
\]
has rational coordinates $e_i$ converging to $\tau$. Sufficiently accurate approximations make every $e_i>0$. This gives the asserted reserve with at most $2k$ atoms. A rational linear dependence among more than $k$ atoms can be eliminated by moving the rational coefficient vector to the first nonnegative boundary along a rational null vector. Each step removes an atom and preserves rational nonnegative coefficients. Lemma~\ref{lem:reserve} then converts weights to Gram columns.
\end{proof}

\begin{lemma}[Facewise stationarity]\label{lem:stationarity}
Let $Q$ be a closed polyhedral vector cone, let $W$ be nonnegative on $Q$, and suppose $x\in Q$ satisfies $x\trans Wx=0$. If $F$ is the minimal face of $Q$ containing $x$, then the restriction of $W$ to $U=\lin F$ is positive semidefinite and annihilates $x$.
\end{lemma}
\begin{proof}
Since $x\in\ri F$, for every $v\in U$ the points $x+tv$ lie in $F$ for sufficiently small positive and negative $t$. The polynomial
\[
 (x+tv)\trans W(x+tv)=2t\,x\trans Wv+t^2v\trans Wv
\]
is nonnegative on a two-sided interval about zero. Its linear coefficient vanishes and its quadratic coefficient is nonnegative. This proves both assertions. The proof includes lineality and $x=0$.
\end{proof}

\begin{lemma}[Exact rational zero domains]\label{lem:zeros}
Let $W$ be a real algebraic symmetric matrix nonnegative on $P$. For every face $F$ of every component, choose a rational basis matrix $B_F$ for $\lin F$. In a common number field $\Q(\theta)$, using an irreducible power basis, expand
\begin{equation}\label{eq:coeff}
 B_F\trans WB_F=\sum_{a=0}^{e-1}\theta^a M_{F,a},\qquad M_{F,a}\in\Q^{r_F\times r_F}.
\end{equation}
Set
\begin{equation}\label{eq:DF}
 D_F=F\cap\left\{B_Fz:M_{F,a}z=0\ \text{for every }a\right\}.
\end{equation}
Then
\[
 \cl\{u\in P\cap\Q^n:u\trans Wu=0\}=\bigcup_F D_F.
\]
In particular this closure is a finite union of closed rational polyhedral cones, contained in the real zero set.
\end{lemma}
\begin{proof}
Take a rational zero $u$ and its minimal component face $F$. Its coordinates $z$ in $B_F$ are rational. Lemma~\ref{lem:stationarity} gives $(B_F\trans WB_F)z=0$. Linear independence of the power basis in \eqref{eq:coeff} gives $M_{F,a}z=0$ for every $a$, so $u\in D_F$.

Conversely every point of $D_F$ has zero quadratic value. Each $D_F$ is rational polyhedral, so its rational points are dense in it. The finite union is closed. The two inclusions now give the identity. No positive-semidefiniteness test for arbitrary $F$ is needed in the definition: stationarity is used only to prove that every rational zero is captured.
\end{proof}

\section{Proof and canonicity of the carrier theorem}\label{sec:carrierproof}
The algorithm maintains a finite union of rational polyhedral cones, initially $P$.

\begin{quote}
\textbf{Arithmetic carrier algorithm.}
Compute $V=\lin C(P)$ and a rational reserve $B$ as in Lemma~\ref{lem:reserve}. Decide whether $C(P)\cap H$ is empty. If empty, return EMPTY. If $H\cap\ri C(P)\ne\varnothing$, return $P$. Otherwise select an algebraic $W$ that is nonnegative on $P$, nonzero on $V$, zero on $H\cap V$, and normalized by $\ip{W}{B}=1$. Replace $P$ by the union in Lemma~\ref{lem:zeros}, and repeat.
\end{quote}
Every real test and selection in this description has an exact real-algebraic implementation; its cost is included in Section~\ref{sec:complexity}.

\begin{lemma}[Availability of an exposer]\label{lem:exposer}
If $C(P)\cap H\ne\varnothing$ but $\ri C(P)\cap H=\varnothing$, the normalized algebraic exposer required by the algorithm exists.
\end{lemma}
\begin{proof}
Work in $V$. Separate the relatively open convex set $\ri C(P)$ from the affine space $H\cap V$. The separating functional is constant on that affine space. Scaling points of the cone and using a feasible point shows that this constant is zero and that the functional is nonnegative on the cone. It is nonzero on $V$. Since $B\in\ri C(P)$, its value at $B$ is strictly positive, so normalization is possible.

The requirements are a first-order real formula over $\Q$: linear equations express vanishing on the rational affine space and the normalization, while a universal vector block expresses nonnegativity on the finite union of rational polyhedral cones. A nonempty semialgebraic set defined over $\Q$ has a point with real algebraic coordinates. Thus an algebraic $W$ can be selected. The zero-dimensional case has already terminated, since feasibility there implies relative strict feasibility.
\end{proof}

\begin{proof}[Proof of Theorem~\ref{thm:main}]
At one reduction let $P'$ be the new domain. It lies in $P$, and $C(P')\subseteq C(P)\cap W^\perp$. Because $W$ is nonzero on $V$, the matrix span strictly decreases.

Every rational feasible Gram factorization survives \emph{column by column}. Indeed, if $X=\sum_j u_ju_j\trans\in H$, then
\[
 0=\ip{W}{X}=\sum_j u_j\trans Wu_j
\]
is a sum of nonnegative quantities. Every $u_j$ is a rational zero and belongs to $P'$ by Lemma~\ref{lem:zeros}. Conversely $P'\subseteq P$. Thus the rational fiber is preserved exactly at each step. There can be at most $d$ proper reductions. A later empty real-feasibility test proves that the original rational fiber was empty.

Suppose the procedure terminates at a nonempty $P_*$. The affine space $H\cap\lin C(P_*)$ is rational, and its intersection with $\ri C(P_*)$ is nonempty and relatively open. It therefore contains a rational point $X_0$. Lemma~\ref{lem:interior} factors $X_0$ rationally and gives a full-span reserve.

For any $X\in C(P_*)\cap H$, the segment $(1-t)X+tX_0$ lies in $H\cap\ri C(P_*)$ for $0<t\le1$. Rational points in the rational affine space can approximate these segment points while remaining in that relatively open set. Lemma~\ref{lem:interior} factors all such rational approximants. This proves \eqref{eq:closure}; the reverse closure inclusion follows from closedness. Exact preservation gives \eqref{eq:exactfiber}.

It remains to prove that $P_*$ itself is intrinsic. If $u\in\mathcal D_{\Q}(P,H)$, decompose both the residual in \eqref{eq:viable} and its positive rational weight into rational squares. The preceding columnwise argument shows that every positive rescaling of $u$ occurring in this factorization survives every step. Hence $u\in P_*$. Conversely let $u\in P_*\cap\Q^n$. For sufficiently small rational $\tau>0$, $X_0-\tau uu\trans$ remains in $\ri C(P_*)$. It is rational, and Lemma~\ref{lem:interior} factors it. Thus $u$ is viable. We have proved
\[
 \mathcal D_{\Q}(P,H)=P_*\cap\Q^n.
\]
Rational points are dense in the finite union $P_*$, which is closed. Therefore $P_*=\cl\mathcal D_{\Q}(P,H)=\AC(P,H)$, independent of the exposers. The width and reserve assertions follow from Lemmas~\ref{lem:reserve} and~\ref{lem:interior}.
\end{proof}

\begin{corollary}[A rank-one arithmetic alternative]\label{cor:insertion}
Let $u\in P\cap\Q^n$. Either $u$ is viable, with a rational insertion certificate \eqref{eq:viable}, or the carrier algorithm produces an empty-fiber certificate or a first exposer strictly positive at $u$. In the latter case the exposer is nonnegative on the \emph{current} domain and annihilates its feasible matrix slice. The preceding reduction chain is part of the certificate.
\end{corollary}
\begin{proof}
A rational zero always survives an update. Thus a rational direction first removed from a nonempty chain has strictly positive value at that step. The converse follows by pairing a hypothetical surviving rational factorization with that exposer.
\end{proof}

A full-span reserve makes insertion quantitative. Write
\[
 X_0=Y+\sum_{i=1}^k e_iG_i,\quad e_i>0,\quad Y\in\GG(\AC),\qquad uu\trans=\sum_{i=1}^ka_iG_i.
\]
Then any rational
\begin{equation}\label{eq:tauinsert}
 0<\tau\le\frac{\min_i e_i}{2(1+\max_i|a_i|)}
\end{equation}
works, since every $e_i-\tau a_i$ remains positive. Cramer's rule bounds the extra binary height polynomially in that of the stored reserve and $u$, for fixed $n$. The same construction inserts any finite list of viable rational directions simultaneously, after reducing their coefficients sufficiently.

